\chapter{Realizácia}
\label{ch:implementation}

Každá konfigurácia z tabuľky~\ref{tab:configurations} je v repozitári súbor manifestov a nastavení Helm~\cite{helm}. Prepína ich jeden skript, ktorý cluster prevedie presne do zvolenej konfigurácie bez ohľadu na predchádzajúci stav. Zmení iba to, čo sa líši, a pody reštartuje len vtedy, keď sa mení niečo, čo Kubernetes nastavuje pri ich vytvorení (pripojenie tokenu, sidecar). Nakoniec počká, kým aplikácia nezačne odpovedať po celej ceste \texttt{productpage}~$\rightarrow$ \texttt{reviews}~$\rightarrow$ \texttt{ratings}~$\rightarrow$ MongoDB, takže test ani meranie nezačne na neustálenom clustri. Tabuľka~\ref{tab:implementation} priraďuje konfiguráciám súbory a príkazy.

\begin{table}[ht!]
  \centering
  \caption{Realizácia konfigurácií v repozitári}
  \label{tab:implementation}
  \footnotesize
  \renewcommand{\arraystretch}{1.15}
  \begin{tabularx}{\textwidth}{@{}>{\raggedright\arraybackslash}p{3.2cm}>{\raggedright\arraybackslash}X>{\raggedright\arraybackslash}p{3.6cm}@{}}
    \toprule
    Konfigurácia & Súbory & Zmena za behu \\
    \midrule
    Bez ochrany & \texttt{infra/cilium-values.yaml} & -- \\
    Segmentácia siete & \texttt{manifests/s1-segmentation/} (10 NetworkPolicy), \texttt{automountServiceAccountToken: false} na každom ServiceAccounte & politiky okamžite; tokeny po reštarte podov \\
    Šifrovanie medzi uzlami & \texttt{manifests/s2a-wireguard/} \texttt{cilium-values.yaml} & reštart agentov Cilium \\
    mTLS medzi službami & \texttt{manifests/s2b-mtls/}, \texttt{infra/istio.sh} & inštalácia Istio, reštart podov so sidecarmi \\
    Mesh bez mTLS & \texttt{manifests/s2b-plain/} (patch \texttt{DISABLE}) & ako mTLS medzi službami \\
    L7 autorizácia & \texttt{manifests/s3-authz/} (6 AuthorizationPolicy) & okamžite cez xDS \\
    \bottomrule
  \end{tabularx}
\end{table}

\section{Cluster a sieť}

Uzly vytvorí \texttt{make up}: tri virtual machines Lima, na nich k3s~\cite{k3s} bez vlastného CNI, vynucovania politík a kube-proxy (prepínače \texttt{-{}-flannel-backend=none}, \texttt{-{}-disable-network-policy}, \texttt{-{}-disable-kube-proxy}), takže sieť podov, vynucovanie politík aj preklad adries služieb zabezpečuje výlučne Cilium. Výpis~\ref{lst:cilium} ukazuje jeho základné nastavenie, spoločné pre všetky konfigurácie.

\begin{lstlisting}[language={},basicstyle=\ttfamily\footnotesize,caption={Základné nastavenie Cilium (výňatok)},label={lst:cilium}]
kubeProxyReplacement: true
routingMode: tunnel
tunnelProtocol: vxlan
socketLB:
  hostNamespaceOnly: true
cni:
  exclusive: false
hubble:
  enabled: true
\end{lstlisting}

Nastavenia \texttt{socketLB.hostNamespaceOnly} a \texttt{cni.exclusive} slúžia na spoluprácu s Istio; druhé ponecháva konfiguráciu iných CNI pluginov (variant Istio CNI) a v režime s init containerom sa neuplatní. Cilium bez prvého z nich prekladá adresu služby na adresu podu už pri volaní \texttt{connect()} v sieťovom namespace podu; sidecar by potom videl spojenie priamo na pod a jeho pravidlá presmerovania by sa minuli účinkom~\cite{cilium-kpr}. Prekladom iba v namespace hostiteľa sa zachová adresa služby, ktorú Envoy potrebuje na smerovanie. Keďže sú tieto nastavenia aj v konfigurácii bez ochrany, konfigurácia mTLS medzi službami nemení nastavenie CNI.

\section{Segmentácia siete}

Segmentácia stojí na štandardných objektoch NetworkPolicy~\cite{k8s-networkpolicy}, ktoré Cilium vynucuje v eBPF podľa identity odvodenej z labelov~\cite{cilium-identity}. Politika \texttt{default-deny} vyberá všetky pody namespace a zakazuje prichádzajúcu aj odchádzajúcu prevádzku; ďalšia politika povoľuje iba DNS. Každý workload má potom vlastnú politiku, ktorá povolí presne jeho hrany v grafe volaní a port. Výpis~\ref{lst:netpol} uvádza politiku pre \texttt{ratings}: prijíma iba od \texttt{reviews} a odchádza iba do MongoDB.

\begin{lstlisting}[language={},basicstyle=\ttfamily\footnotesize,caption={NetworkPolicy pre \texttt{ratings}},label={lst:netpol}]
spec:
  podSelector:
    matchLabels: {app: ratings}
  policyTypes: [Ingress, Egress]
  ingress:
    - from:
        - podSelector:
            matchLabels: {app: reviews}
      ports:
        - {protocol: TCP, port: 9080}
  egress:
    - to:
        - podSelector:
            matchLabels: {app: mongodb}
      ports:
        - {protocol: TCP, port: 27017}
\end{lstlisting}

Okrem piatich politík aplikácie a dvoch základných sú v sade tri politiky pre pody merania, aby záťažové nástroje nepotrebovali výnimku z politiky default-deny. V konfiguráciách s mesh pribudne politika, ktorá všetkým podom povolí spojenie na istiod (port 15012); bez nej by sidecary nedostali konfiguráciu ani certifikát. Segmentácia teda musí poznať aj riadiacu rovinu mesh.

Hardening tokenov nastaví \texttt{automountServiceAccountToken: false} na každom ServiceAccounte namespace~\cite{k8s-serviceaccounts}. Pod tak nedostane token, ktorým by sa mohol prihlásiť k API serveru; aplikácia Bookinfo API server nepoužíva, takže jej to nevadí. Každý workload, vrátane MongoDB, má vlastný ServiceAccount, aby mal v mesh vlastnú identitu.

\section{Šifrovanie medzi uzlami}

Transparentné šifrovanie zapína v Cilium dvojica nastavení \texttt{encryption.enabled: true} a \texttt{encryption.type: wireguard}~\cite{cilium-wireguard}. Agent Cilium na každom uzle vygeneruje pár kľúčov, verejný kľúč zverejní v objekte CiliumNode a s ostatnými uzlami vytvorí rozhranie \texttt{cilium\_wg0}. Prevádzka podov smerujúca na iný uzol potom prechádza týmto rozhraním v jadre ako UDP na porte 51871~\cite{donenfeld2017wireguard}. Prevádzka hostiteľov (kubelet, API server) ostáva nešifrovaná, lebo útok odpočúvaním sa týka prevádzky podov. Zmena sa prejaví po reštarte agentov Cilium a aplikácia ani pody o nej nevedia. Vedľajším účinkom je menšie MTU trás podov, ktoré ovplyvní aj prevádzku na jednom uzle (kapitola~\ref{ch:overhead}).

\section{mTLS medzi službami}

Istio sa inštaluje cez Helm (charty \texttt{base} a \texttt{istiod}) a istiod beží na riadiacom uzle, aby nesúperil s aplikáciou o CPU pracovných uzlov. Label \texttt{istio-injection=enabled} na namespace zapne vkladanie sidecarov; presmerovanie prevádzky do Envoy nastaví init container pravidlami iptables~\cite{istio-security}. Nastavenie \texttt{holdApplicationUntilProxyStarts} spustí aplikáciu až po pripravenom sidecari, takže prvé spojenia nezlyhajú. Istio vydá každému podu certifikát s identitou v tvare SPIFFE, napríklad \url{spiffe://cluster.local/ns/bookinfo/sa/bookinfo-reviews}~\cite{spiffe-id}. Vynucovanie mTLS zapína PeerAuthentication na úrovni namespace (výpis~\ref{lst:peerauth})~\cite{istio-peerauthentication}.

\begin{lstlisting}[language={},basicstyle=\ttfamily\footnotesize,caption={PeerAuthentication pre namespace \texttt{bookinfo}},label={lst:peerauth}]
apiVersion: security.istio.io/v1
kind: PeerAuthentication
metadata:
  name: default
spec:
  mtls:
    mode: STRICT
\end{lstlisting}

Mesh bez mTLS používa tie isté manifesty s režimom \texttt{DISABLE}. Pre istotu táto konfigurácia pridáva aj DestinationRule, ktorá šifrovanie výslovne vypne aj na strane klienta, takže výsledok nezávisí od automatického mTLS.

\section{L7 autorizácia}

Autorizácia používa vzor \enquote{allow-nothing}: AuthorizationPolicy bez pravidiel zakáže každú požiadavku na workload v namespace a pre každý workload pribudne jedna politika \texttt{ALLOW}, ktorá povolí iba jeho volajúcich podľa identity z mTLS, metódy a cesty~\cite{istio-authorizationpolicy}. Výpis~\ref{lst:authz} uvádza politiku pre \texttt{ratings}: \texttt{reviews} smie hodnotenia iba čítať.

\begin{lstlisting}[language={},basicstyle=\ttfamily\footnotesize,caption={AuthorizationPolicy pre \texttt{ratings}},label={lst:authz}]
spec:
  selector:
    matchLabels: {app: ratings}
  action: ALLOW
  rules:
    - from:
        - source:
            principals:
              - cluster.local/ns/bookinfo/sa/bookinfo-reviews
      to:
        - operation:
            methods: [GET]
            paths: [/ratings/*]
\end{lstlisting}

MongoDB komunikuje vlastným protokolom nad TCP, preto jej politika rozlišuje iba identitu \texttt{ratings} a port. Vstupný bod \texttt{productpage} prijíma požiadavky od kohokoľvek, ale iba metódy na čítanie a prihlasovací formulár. Politiky neovplyvnia kontroly stavu od kubeletu, lebo tie injektor presmeruje na agenta sidecaru. Politiky sa do sidecarov dostanú cez xDS za niekoľko sekúnd a pody sa nereštartujú.

\section{Automatizácia}

Všetky kroky sú ciele súboru \texttt{Makefile}: \texttt{make up} vytvorí cluster, \texttt{make app} nasadí aplikáciu, \texttt{make s0} až \texttt{make s3} a \texttt{make s2b-plain} prepnú konfiguráciu, \texttt{make test} spustí všetky útoky proti aktuálnej konfigurácii, \texttt{make bench} meranie réžie a \texttt{make plot} z nameraných dát vytvorí tabuľky a grafy (knižnica Matplotlib~\cite{matplotlib}). Každý test zapíše riadok do výsledkového CSV s identifikátorom behu a commitom repozitára a uloží dôkaz (výstup príkazov, záznam prevádzky) pod názvom, ktorý určuje útok aj konfiguráciu. Verzie všetkých komponentov sú pripnuté a pri každom behu sa zaznamenávajú.
